\scanpage{11}
\mspage{4}
$\mathrm{Hom}$ internes existent dans \deleted{les} $\mathrm{Ho}_W$ et les $\mathrm{Ho}_{W(I)}$, et se calculent comme dit \unsure{dans|sur} la prop.

Le seul problème est l'existence des $\underline{\mathrm{Hom}}$ internes dans les $M(I)$, on aura
\[
\begin{aligned}
\underline{\mathrm{Hom}}_{M(I)}(X,Y)
 &=\bigl(i\longmapsto
 \underbrace{\underset{\substack{\downarrow\\\inserted{\substack{\text{\unsure{à val.} dans}\\M}}}}{\underline{\mathrm{Hom}\,\mathrm{ext}}}
 (X\vert I_{/i},Y\vert I_{/i})}_{\displaystyle
 \subset\prod_{\substack{\deleted{\unsure{(u,j)\in}}\\(\unsure{j,u})\in\mathrm{Ob}\,I_{/i}}}\underline{\mathrm{Hom}}_M(X(j),Y(j))}
\end{aligned}
\]

- - - -

\uline{NB} \inserted{$(X,Y)\mapsto$} $\underline{\mathrm{Hom}}(X,Y)$ définit un morphisme de dérivateurs
\[
\underset{\diagup}{\mathbb{D}^{\circ}}\times\underset{\diagup}{\mathbb{D}}\unsure{-|\to}\underset{\diagup}{\mathbb{D}}
\]
Si $f:I'-I$
\[
 f^*(\underline{\mathrm{Hom}}_{\mathbb{D},I}
 (\underset{\substack{\uparrow\\\inserted{\mathbb{D}(I)}}}{X},\underset{\uparrow}{Y}))
 \xrightarrow{\sim}\underline{\mathrm{Hom}}_{\mathbb{D},I'}(f^*X,f^*Y)
\]
\margin{Non !}

\uline{Plutôt} Si $X_{.}\in\mathbb{D}(I)$, $Y_{.}\in\mathbb{D}(J)$, alors on définit
\[
\begin{aligned}
\underline{\mathrm{Hom}}_{\square}(X_{.},Y_{.})&\in\mathbb{D}(I^\circ\times J)=\mathbb{D}_{I^\circ}(J)\\
&=\mathbb{D}_J(I^\circ)
\end{aligned}
\]
Dont il faudrait donner une car. intrinsèque en termes de $\mathbb{D}$. \deleted{\illegible{1}}
% archive sheet number: 10

